% Hand-written: staged design rationale (stage / role / benefit / limitation)
\begin{table}[t]
\centering
\small
\setlength{\tabcolsep}{4pt}
\begin{tabular}{@{}p{1.4cm}p{2.2cm}p{1.7cm}p{1.7cm}@{}}
\toprule
\textbf{Stage} & \textbf{Role} & \textbf{Benefit} & \textbf{Limitation} \\
\midrule
Gemma extraction & resolve obvious (person) synonyms while reading & typed, span-grounded attributes; free bulk resolution & selective, inconsistent, can mis-extract \\
\addlinespace
v1 (round 1) & merge frequent residual variants; blacklist noise & cheap; completes the head; removes worst noise & misses long tail / phrases \\
\addlinespace
v2 (converged, judge feedback) & merge long-tail and phrasal variants; remove judge-flagged noise & reaches rare/phrasal synonyms; targeted cleanup & not exhaustive; needs a judge; ongoing \\
\bottomrule
\end{tabular}
\caption{The three synonym-handling stages are complementary. Extraction resolves the
obvious synonyms, v1 the frequent residual, v2 the long tail; each recovers cases the
others cannot. The loop relabels and filters at the type level and never re-verifies an
individual extraction, so it is precision-preserving: its gains are recall (via merging)
and facet usability (via filtering), and each step can be run by a human or by the agent.}
\label{tab:design-rationale}
\end{table}
